\section{Discussion}
\label{sec:discussion}

\subsection{Why fewer physically entitled pages can be faster}

\system does not make SSD reads free and does not eliminate expert misses. Its
Qwen hit rate remains near 87\%, below fixed-96 near 91\% and far above
fixed-16 near 49\%. The improvement comes from a different system balance. The
large logical cache still captures reuse when cold pages survive, while the OS
can reclaim cold pages instead of compressing them or displacing more useful
file-backed state. The fixed-96 baseline shows that maximizing hit rate can
reduce service rate once its physical claim harms the shared pool.

This interpretation is supported by three observations. First, physical
footprint and compressor growth fall materially relative to fixed-96. Second,
the candidate leads in no-pressure experiments as well as under the controlled
co-tenant, so the result is not solely an artifact of one pressure process.
Third, repeated empty recovery occurs in every candidate row without output
divergence. The mechanism is therefore exercised in the measured path rather
than acting as an unused safety annotation.

\subsection{The OS as eviction partner, not model scheduler}

A recent kernel-cache study finds that ordinary OS replacement can be a strong
expert tier when the kernel sees the complete physical workload
\cite{si2026kernelcache}. \system reaches a related design point without
delegating model identity or exact publication to the kernel. The runtime
decides which pages are eligible, retains hot experts, validates slot contents,
and reloads exact weights. macOS decides whether eligible pages should remain
physical given global demand. This division avoids a runtime-only estimate of
desktop competition while preserving model-specific correctness.

The design is not a complete adaptive memory controller. It uses a fixed
logical/hot split and lets the OS react within the cold set. That simplicity is
a strength of the current evidence: there is no learned predictor, online
exploration, or second model configuration that can perturb the state being
measured. Future work can adapt the hot-set size, but it must beat this
non-destructive baseline and account for transition cost.

\subsection{Relationship to discrete-GPU edge serving}

FreeToken treats CPU memory, GPU memory, PCIe bandwidth, and CPU compute as
heterogeneous resources and continuously maps work across them
\cite{yang2026freetoken}. \system adopts its challenge-driven systems
discipline but not its CUDA data plane. On Apple Silicon there is no expert copy
across a host/device boundary. The relevant elasticity is whether cold shared
pages remain physically valid. CPU execution, bandwidth-adaptive miss splitting,
semantic checkpoints, and agent lifecycle management remain complementary.

\subsection{Implications for model support}

The mechanism transfers from Qwen's 256-expert, 40-layer layout to Gemma's
128-expert, 30-layer layout without policy tuning. That supports an
architecture-independent cache-validity abstraction, not arbitrary-model
compatibility. A new model still needs a correct tokenizer, packed manifest,
tensor mapping, router semantics, quantization decoder, and Metal kernels.
\system's public catalog makes checkpoint metadata data-driven, while the
native runtime remains the correctness boundary.

\subsection{Practical deployment}

The public CLI defaults to OS-managed 96/16 residency for admitted profiles and
offers one-shot generation plus a loopback OpenAI/Anthropic-shaped server. It
reuses one canonical model store, pins revisions, serializes downloads, and
refuses a transfer that would violate its disk reserve. These product surfaces
are engineering contributions for reproducibility and adoption; they are not
part of the measured algorithmic claim.
